\subsection{Sparse variational inference}
\label{sec:variational}

Exact inference in the model of Eq.~\eqref{eq:model} is impractical for two reasons: the covariance matrix of the Gaussian process grows quadratically with the number $N$ of cell--day observations and its factorization cubically, and the hurdle likelihood is not conjugate to the Gaussian prior, so the posterior has no closed form.
We therefore approximate the latent field with $M=300$ inducing variables $\mathbf{u}$ at learnable space--time positions $\mathbf{Z}$~\citep{titsias2009variational,hensman2015scalable}, with prior and variational posterior
\begin{equation}
    p(\mathbf{u})=\mathcal{N}(\mathbf{0},K_{uu}),
    \qquad
    q(\mathbf{u})=\mathcal{N}(\mathbf{m},S).
\end{equation}
At a set of space--time locations $R$, the latent field is the projection
\begin{equation}
    \mathbf{f}_R=A_R\mathbf{u},
    \qquad
    A_R=K_{Ru}K_{uu}^{-1},
\end{equation}
so that $q(\mathbf{f}_R)=\mathcal{N}(A_R\mathbf{m},A_RSA_R^\top)$. This is the subset-of-regressors approximation of the Gaussian process~\citep{quinonero2005unifying}, a rank-$M$ approximation that, unlike a standard sparse variational Gaussian process, omits the residual covariance $K_{RR}-A_RK_{uR}$.

The parameters are learned by maximizing the evidence lower bound
\begin{equation}
    \mathcal{L}
    =\sum_{i=1}^{N}\mathbb{E}_{q(\mathbf{u})}\bigl[\log p(y_i\mid\mathbf{u},\boldsymbol{\phi})\bigr]
    -\operatorname{KL}\bigl(q(\mathbf{u})\,\|\,p(\mathbf{u})\bigr),
\end{equation}
where $p(y_i\mid\mathbf{u},\boldsymbol{\phi})$ is the hurdle likelihood of Eq.~\eqref{eq:hurdle}. With $b_i=\mathbf{1}[y_i>0]$, its logarithm splits into two terms,
\begin{equation}
    \begin{aligned}
        \log p(y_i\mid\mathbf{u},\boldsymbol{\phi})
        &=b_i\log p_i+(1-b_i)\log(1-p_i)\\
        &\quad+b_i\log\frac{\operatorname{Poisson}(y_i\mid\lambda_i)}{1-\operatorname{Poisson}(0\mid\lambda_i)}.
    \end{aligned}
\end{equation}
The first is the Bernoulli log-likelihood of activity, which depends only on the classifier, and the second is the zero-truncated Poisson log-likelihood of the positive counts, which depends only on the LGCP. The classifier therefore learns from all cell--days whether activity occurs, and the LGCP from the active cell--days how much activity there is.
The KL term is evaluated analytically, whereas the expected zero-truncated Poisson term is estimated by Monte Carlo sampling of $\mathbf{u}$ on minibatches of cell--days. We maximize $\mathcal{L}$, with an $L_2$ penalty on the classifier weights, using Adam jointly over the variational parameters $\mathbf{m}$ and $S$, the inducing positions $\mathbf{Z}$, the coefficients $\alpha$ and $\boldsymbol{\beta}$, the kernel variances and period, and the classifier weights $\boldsymbol{\phi}$. At prediction time, all expectations under $q$ are estimated with 300 samples of the latent field.
